% =============================================================================
% SULI 2026 — ML K+ PID for CLAS12 SIDIS
% Author: Cooper Bell (SULI Student, Jefferson Lab)
% PI: Maria Zurek (Argonne National Laboratory)
%
% TEMPLATE — fill in content as the summer progresses.
% Weeks 1-2 deliverable: Sections 1-5 partially filled; placeholders
% for all figures and tables Cooper produces in Weeks 1-2.
%
% Compile: pdflatex report_template.tex  (twice for cross-references)
%          bibtex report_template
%          pdflatex report_template.tex (twice more)
% =============================================================================

\documentclass[11pt]{article}

\usepackage[margin=1in]{geometry}
\usepackage{graphicx}
\usepackage{amsmath}
\usepackage{amssymb}
\usepackage{hyperref}
\usepackage{siunitx}
\usepackage{booktabs}
\usepackage{cite}

\hypersetup{colorlinks=true, linkcolor=blue, citecolor=blue, urlcolor=blue}

% -----------------------------------------------------------------------------
% Handy shorthand — add sparingly; keep the preamble minimal
% -----------------------------------------------------------------------------
\newcommand{\Kp}{K^{+}}
\newcommand{\pip}{\pi^{+}}
\newcommand{\epKX}{ep \to e'\,p\,\Kp\,X}

% =============================================================================
\title{Machine-Learning Particle Identification\\
       for \boldmath{$K^+$} Identification in CLAS12 SIDIS}

\author{Cooper Bell \\
  \small{SULI Student, Jefferson Lab} \\[2pt]
  \small{PI: Maria Zurek, Argonne National Laboratory}}

\date{\today}
% =============================================================================

\begin{document}

\maketitle

% =============================================================================
\begin{abstract}
% Cooper: write ~150 words here by end of Week 2. Describe (1) the physics
% motivation — kaon SIDIS and why PID purity matters; (2) the contamination
% problem at high momentum; (3) the ML approach (BDT trained on clasdis MC);
% (4) the scope of this report (Weeks 1-2: ntuple, baseline, data/MC audit).
% Draft placeholder below — replace completely.
TODO: Abstract placeholder. This report describes the development of a
machine-learning classifier for $\Kp$ particle identification in the
CLAS12 SIDIS reaction $\epKX$.
The study uses \texttt{clasdis} Monte Carlo with geometric truth matching
to characterize baseline Event Builder PID performance and to audit detector
variable agreement between simulation and RGA pass-2 data.
Weeks~1 and~2 deliverables are documented here; model training and
validation will be reported in subsequent updates.
\end{abstract}

% =============================================================================
\section{Introduction}
\label{sec:intro}
% Cooper: 3-4 paragraphs. Cover:
%   (1) Why kaon SIDIS matters — flavor decomposition of TMDs, strangeness.
%   (2) The contamination problem: pion leakage into EB-K+ sample at high p.
%       Quote the rough numbers (30-50% at p > 2.5 GeV/c) from Week-2 baseline.
%   (3) Why ML is a natural fit: chi2pid uses one FTOF number; ML can use all
%       detector responses simultaneously (FTOF panels, ECAL, HTCC).
%   (4) State the analysis channel and Cooper's specific goal.
% Reference the CLAS12 detector paper \cite{Burkert2020} and FTOF paper
% \cite{Carman2020} where appropriate.
% -----------------------------------------------------------------------

Semi-inclusive deep-inelastic scattering (SIDIS) off proton targets
provides direct access to transverse-momentum-dependent parton distributions
(TMDs) and their flavor structure.
Kaon electroproduction, $\epKX$, is particularly sensitive to the
strange-quark sea and to fragmentation functions that are not constrained by
pion measurements alone.
Accurate $\Kp$ identification is therefore essential for the SIDIS physics
program at CLAS12~\cite{Burkert2020}.

The CLAS12 Event Builder (EB) assigns particle identities using time-of-flight
information from the FTOF system~\cite{Carman2020}.
At low momenta ($p \lesssim \SI{2}{GeV/c}$), pions and kaons have
sufficiently different velocities ($\beta = v/c$) that the FTOF resolves them
cleanly.
Above $p \approx \SI{2.5}{GeV/c}$ the $\beta$ curves for $\pip$ and $\Kp$
converge, and pion contamination in the EB-identified $\Kp$ sample can reach
30--50\% (to be quantified in Section~\ref{sec:baseline}).
The EB's refinement quantity \texttt{chi2pid} encodes the goodness-of-fit of
the measured $\beta$ to the kaon hypothesis; the production pass-2 cut is
momentum-dependent, but still uses essentially one number from the FTOF.

A machine-learning classifier can exploit the full set of detector responses
available per track---FTOF hit patterns across two panel layers, calorimeter
energy deposits (PCAL, ECin, ECout), and the HTCC Cherenkov signal---to learn
discriminants that \texttt{chi2pid} alone misses.
The primary model is a gradient-boosted decision tree (BDT) trained on
\texttt{clasdis} Monte Carlo, where geometric truth matching provides exact
particle labels.

This report documents Weeks~1 and~2 of the 10-week project:
ntuple production, baseline PID characterization, and the data/MC variable
agreement audit that governs which features enter ML training.

% =============================================================================
\section{Data and Monte Carlo Samples}
\label{sec:samples}

\subsection{Monte Carlo: \texttt{clasdis}}
% Cooper: describe clasdis briefly — what generator it is, inbending RGA
% kinematics, beam energy, approximate file count you processed.
% Note the truth-matching algorithm (anchor on MC::Particle, geometric cone
% |Delta phi| < 6 deg, |Delta theta| < 2 deg, best match wins).
The \texttt{clasdis} generator~\cite{Burkert2020} provides the Monte Carlo
sample for this study.
We use the RGA fall-2018 inbending configuration at a beam energy of
\SI{10.604}{GeV}.
% TODO: insert total file count and approximate event yield after Week-2 scale-up.
Each reconstructed track is geometrically matched to a \texttt{MC::Lund} truth
particle within a cone of $|\Delta\phi| < 6^\circ$ and
$|\Delta\theta| < 2^\circ$; the closest match within the cone is selected.
Tracks with no match are dropped from the training set.

\subsection{RGA Pass-2 Data}
% Cooper: briefly describe the RGA pass-2 data subset you compare against in
% the data/MC audit (Section 5). Note FD-only, EB-K+ tracks, approximate
% run range if known. Exact run list TBD with Maria.
RGA pass-2 data from the fall-2018 inbending run period are used for the
data/MC variable agreement study (Section~\ref{sec:datamc}).
Only EB-identified $\Kp$ tracks in the Forward Detector are selected,
applying the same vertex and fiducial requirements as the MC sample.
The run list and luminosity normalization are confirmed with the PI.

\subsection{Ntuple Production Pipeline}
\label{sec:pipeline}
% Cooper: describe the processing chain in your own words after running it.
% Mention both groovy scripts, the converter, and the 53-column output.
% Describe the cuts explicitly (FD status in [2000,4000), vertex |vz| range,
% DC fiducial flags). Confirm exact column count once you have run the script.

Ntuples are produced from HIPO files using the CLAS12 analysis software
framework~\cite{Burkert2020}.
Two Groovy scripts are used:
\texttt{processing\_mc\_pid\_training.groovy} processes MC files and writes
MC-truth labels; \texttt{processing\_data\_pid\_training.groovy} runs the
identical reconstruction chain on data without truth information.
Both produce space-separated text files that are converted to ROOT TTrees
by \texttt{convert\_txt\_to\_root.cpp}.

\textbf{Selection cuts applied at the ntuple stage:}
\begin{itemize}
  \item Forward Detector only: track \texttt{status} $\in [2000, 4000)$.
  \item Vertex $z$-position: $|v_z|<$ [TODO: confirm cut value with Maria].
  \item DC fiducial requirements: standard CLAS12 fiducial flags.
  \item EB-assigned PID restricted to $\{\pip, \pi^-, \Kp, K^-\}$.
\end{itemize}

Each output row corresponds to one reconstructed hadron track and contains
53~columns organized as follows:
event identifiers (run, event, helicity);
DIS kinematics ($Q^2$, $W$, $x$, $y$, $\nu$);
per-track kinematics and EB identity ($p$, $\theta$, $\phi$, $v_z$, sector,
status, \texttt{pid});
raw EB inputs (\texttt{beta}, \texttt{chi2pid});
FTOF panel-1A/1B responses (energy, time, path length);
ECAL inner/outer responses (energy, time, path length);
HTCC photoelectron count (\texttt{nphe\_htcc});
PCAL and FTOF panel-2 responses (energy, time, path length);
14~RICH variables (cross-check only, not used in ML training);
and MC truth columns (\texttt{mc\_matching\_pid}, \texttt{mc\_parent\_pid},
\texttt{mc\_match\_quality}).
Missing values use the sentinel $-9999$.
The full column reference is given in Appendix~\ref{app:columns}.

% =============================================================================
\section{PID Metric Definitions}
\label{sec:metrics}
% Cooper: this section is already drafted below using Marco's definitions from
% the 2026-04-21 SIDIS meeting (Mirazita \cite{Mirazita2026}).
% Review it, verify the notation matches what you use in your Python scripts,
% and add one sentence connecting each metric to the pd.crosstab output.

We follow the PID metric conventions established by the CLAS12 SIDIS
group~\cite{Mirazita2026}.
Let $N_{i \to j}$ denote the number of tracks whose MC truth species is $i$
and whose Event Builder (or classifier) assigned species is $j$, where
$i, j \in \{\pip, \Kp, p, \ldots\}$.
All four metrics below are computed from the single confusion matrix
$N_{i \to j}$, obtained in Python via
\texttt{pd.crosstab(df["mc\_matching\_pid"], df["pid"])}.

\textbf{Purity} of the $\Kp$ sample --- fraction of EB-identified $\Kp$ that
are truly $\Kp$:
\[
  P^{K} = \frac{N_{K \to K}}{\displaystyle\sum_{i} N_{i \to K}}
\]

\textbf{Contamination} of $\Kp$ by species $i$ --- fraction of EB-identified
$\Kp$ that are actually species $i$:
\[
  C^{i \to K} = \frac{N_{i \to K}}{\displaystyle\sum_{k} N_{k \to K}}
\]

Closure: $P^{K} + \sum_{i \neq K} C^{i \to K} = 1$.

\textbf{Efficiency} --- fraction of true $\Kp$ correctly identified as $\Kp$:
\[
  \varepsilon^{K} = \frac{N_{K \to K}}{\displaystyle\sum_{j} N_{K \to j}}
\]

\textbf{Mis-ID} of true $\Kp$ to species $j$ --- fraction of true $\Kp$ wrongly
assigned to $j$:
\[
  M^{K \to j} = \frac{N_{K \to j}}{\displaystyle\sum_{k} N_{K \to k}}
\]

Closure: $\varepsilon^{K} + \sum_{j \neq K} M^{K \to j} = 1$.

Purity and contamination normalize to \emph{detected} (EB-labeled) counts and
depend on the $\pip{:}\Kp{:}p$ production-rate ratio in the sample; use them
when reporting the composition of a selected data sample.
Efficiency and mis-ID normalize to \emph{true} counts and reflect detector
and classifier performance independently of the production rate; use them when
characterizing the PID system itself.

% =============================================================================
\section{Baseline PID Performance}
\label{sec:baseline}

\subsection{Standard EB \texttt{chi2pid} Cuts}
% Cooper: describe the three chi2pid options clearly:
%   (1) flat |chi2pid| < 3 (loose, orientation only — not the production baseline)
%   (2) momentum-dependent pass-2 cut (baseline_chi2pid.py: passes_kplus_chi2pid_cut)
%   (3) per-run Sp19/MC variant (if relevant for your sample)
% State clearly which one is the headline baseline and why.

The CLAS12 Event Builder provides a variable \texttt{chi2pid} that quantifies
how closely a track's measured $\beta$ matches the kaon hypothesis given its
momentum.
Three common cut choices appear in the literature:
\begin{enumerate}
  \item \textbf{Flat} $|\texttt{chi2pid}| < 3$: symmetric window used for
        orientation and quick checks; \emph{not} the production baseline.
  \item \textbf{Pass-2 momentum-dependent cut}
        (\texttt{passes\_kplus\_chi2pid\_cut} in
        \texttt{scripts/baseline\_chi2pid.py}): asymmetric window that
        contracts above $p = \SI{2}{GeV/c}$ to compensate for worsening
        $\pip$/$\Kp$ separation.
        This is the headline baseline used throughout this study.
  \item \textbf{Per-run Sp19/MC variant}: an alternative parameterization
        from the Spring-2019 data period; included for comparison where noted.
\end{enumerate}

% Cooper: after computing baseline results in Week 2, add 1-2 sentences
% interpreting the momentum-dependent cut's behavior and why it was chosen.

\subsection{Contamination Matrix from MC Truth}
\label{sec:contam_matrix}
% Cooper: replace the placeholder values in Table 1 with real numbers from
% pd.crosstab(df["mc_matching_pid"], df["pid"]) on the EB-K+ subsample.
% Add statistical uncertainties (sqrt(N) for raw counts; propagate to ratios).

Table~\ref{tab:contam_matrix} shows the EB-pid $\times$ MC-truth contamination
matrix for the \texttt{clasdis} sample after the pass-2 momentum-dependent
\texttt{chi2pid} cut.
Each entry $N_{i \to K}$ is the raw track count; derived metrics
$P^K$ and $C^{i \to K}$ appear in the bottom rows.

\begin{table}[h]
  \centering
  \caption{EB-pid $\times$ MC-truth confusion matrix for EB-identified $\Kp$
           tracks after the pass-2 \texttt{chi2pid} cut.
           Values are placeholder (TBD) --- replace with Week-2 results.}
  \label{tab:contam_matrix}
  \begin{tabular}{lcccc}
    \toprule
    & \multicolumn{4}{c}{MC truth species} \\
    \cmidrule(lr){2-5}
    Quantity & $\Kp$ & $\pip$ & $p$ & other \\
    \midrule
    $N_{i \to K}$ (raw counts)          & \textit{TBD} & \textit{TBD} & \textit{TBD} & \textit{TBD} \\
    $C^{i \to K}$ (contamination)       & ---          & \textit{TBD} & \textit{TBD} & \textit{TBD} \\
    \midrule
    $P^K$ (purity) $= N_{K \to K}/\sum_i N_{i \to K}$ & \multicolumn{4}{c}{\textit{TBD}} \\
    $\varepsilon^K$ (efficiency)        & \multicolumn{4}{c}{\textit{TBD}} \\
    \bottomrule
  \end{tabular}
\end{table}

\begin{figure}[h]
  \centering
  % Cooper: replace placeholder.png with figures/beta_vs_p_truth_classes.png
  % once produced in Week 1 Task 4.
  \includegraphics[width=0.85\textwidth]{placeholder.png}
  \caption{TODO: $\beta$ vs $p$ for EB-K+ tracks, colored by MC truth class
           ($\pip$, $\Kp$, $p$).
           The convergence of the $\pip$ and $\Kp$ bands above
           $p \approx \SI{2.5}{GeV/c}$ illustrates the contamination problem.
           Produce in Week~1 Task~4 and save as
           \texttt{figures/beta\_vs\_p\_truth\_classes.png}.}
  \label{fig:beta_vs_p}
\end{figure}

\begin{figure}[h]
  \centering
  \includegraphics[width=0.85\textwidth]{placeholder.png}
  \caption{TODO: \texttt{chi2pid} distribution per MC truth class ($\pip$,
           $\Kp$, $p$) for all EB-K+ tracks (no \texttt{chi2pid} cut applied).
           The degree of overlap between the $\pip$ and $\Kp$ distributions
           sets the ceiling on \texttt{chi2pid}-based separation.
           Produce in Week~1 Task~4; save as
           \texttt{figures/chi2pid\_by\_truth\_class.png}.}
  \label{fig:chi2pid_dist}
\end{figure}

\begin{figure}[h]
  \centering
  \includegraphics[width=\textwidth]{placeholder.png}
  \caption{TODO: \texttt{chi2pid} vs $p$ 2D maps for EB-K+ tracks, one panel
           per MC truth class.
           Indicates where the pass-2 momentum-dependent cut is tight vs.\
           loose.
           Produce in Week~2 Step~2b; save as
           \texttt{figures/chi2pid\_vs\_p\_2d.png}.}
  \label{fig:chi2pid_2d}
\end{figure}

\subsection{Baseline Efficiency, Purity, and Contamination per $(p,\,\theta)$ Bin}
% Cooper: after running compute_baseline.py (Week-2 Step 2b), replace the
% figure placeholder with the nine-panel 2D map figure
% (figures/baseline_2d_maps.png). Add a paragraph describing the main trends:
% where is contamination worst, where is efficiency highest, etc.

Figure~\ref{fig:baseline_2d} shows the full set of baseline performance metrics
in $(p,\,\theta)$ bins.
% TODO: add 1-2 sentences interpreting the 2D maps after Week-2 results are in.

\begin{figure}[h]
  \centering
  \includegraphics[width=\textwidth]{placeholder.png}
  \caption{TODO: nine-panel 2D map in $(p,\,\theta)$ for the baseline
           pass-2 \texttt{chi2pid} cut.
           Top row: raw counts $N_{K \to K}$, $N_{\pi \to K}$, $N_{p \to K}$.
           Middle row: $P^K$, $C^{\pi \to K}$, $C^{p \to K}$.
           Bottom row: $\varepsilon^K$, $M^{K \to \pi}$, $M^{K \to p}$.
           Produce in Week~2 Step~2b; save as
           \texttt{figures/baseline\_2d\_maps.png}.}
  \label{fig:baseline_2d}
\end{figure}

% =============================================================================
\section{Data/MC Variable Agreement Study}
\label{sec:datamc}

\subsection{Methodology}
% Cooper: describe the audit procedure: coarse (p,theta) bins (3x3 grid),
% KS test threshold (> 0.05 flags a variable), KEEP/CANDIDATE/DROP decision
% logic, and the standing rule that any disagreeing variable is excluded from
% ML training before training begins.

Before any ML training, each detector variable is compared between the
\texttt{clasdis} MC (EB-$\Kp$ tracks) and RGA pass-2 data (EB-$\Kp$ tracks)
in coarse $(p,\,\theta)$ slices: $p \in [1,2],\,[2,3],\,[3,5]$~\GeV{} crossed
with $\theta \in [5^\circ, 15^\circ],\,[15^\circ, 25^\circ],\,[25^\circ,
35^\circ]$, giving nine cells.
In each cell, the 1D distributions of MC and data are overlaid.
A variable is flagged if the Kolmogorov--Smirnov (KS) statistic exceeds 0.05
or if the shapes disagree visually.

Flagged variables are assigned one of three statuses:
\textbf{KEEP} (agreement sufficient for ML training),
\textbf{CANDIDATE} (borderline; requires further evaluation),
or \textbf{DROP} (excluded from ML training until the disagreement is understood).
No flagged variable enters the training feature set until this audit is resolved;
this rule is non-negotiable.

\subsection{Per-Feature Comparison}
% Cooper: replace the placeholder figure with the feature audit overlay grid
% once produced in Week-2 Step 2c. Describe the most striking agreements and
% disagreements in a brief paragraph.

\begin{figure}[h]
  \centering
  \includegraphics[width=\textwidth]{placeholder.png}
  \caption{TODO: data/MC distribution overlays for key ML training features
           in the $(p,\,\theta)$ coarse bins.
           Each panel shows MC (filled) and data (open histogram) for one
           variable in one kinematic cell.
           Features shown: \texttt{beta}, \texttt{chi2pid}, FTOF and ECAL
           energies, \texttt{nphe\_htcc}.
           Save individual panels to
           \texttt{figures/feature\_audit/} and a combined grid as
           \texttt{figures/feature\_audit\_grid.png}.}
  \label{fig:mcdata_audit_ftof}
\end{figure}

\begin{figure}[h]
  \centering
  \includegraphics[width=0.75\textwidth]{placeholder.png}
  \caption{TODO: 2D $(p,\,\theta)$ data/MC ratio map for EB-$\Kp$ tracks.
           A flat ratio map ($\approx 1$ everywhere) indicates no kinematic
           reweighting is needed for ML training; significant variation warrants
           reweighting in Week~4.
           Save as \texttt{figures/feature\_audit/ptheta\_data\_mc\_ratio.png}.}
  \label{fig:ptheta_ratio}
\end{figure}

\subsection{Discussion of Disagreements}
% Cooper: after completing the audit (Week-2 Step 2c/2d), write 1-2 paragraphs
% summarizing: which variables passed (KEEP), which are borderline (CANDIDATE),
% which are excluded (DROP). State whether the (p,theta) ratio map in
% Figure \ref{fig:ptheta_ratio} warrants reweighting. Document the final
% feature list agreed on with Maria.

Table~\ref{tab:feature_audit} summarizes the per-feature KEEP/CANDIDATE/DROP
decisions from the Week-2 audit.

\begin{table}[h]
  \centering
  \caption{Per-feature data/MC agreement audit results (Weeks~1--2).
           KS$>0.05$ or visual shape disagreement triggers a CANDIDATE or DROP
           assignment.
           Values TBD --- fill after completing the audit in Week~2.}
  \label{tab:feature_audit}
  \begin{tabular}{llcc}
    \toprule
    Variable & Group & KS (worst cell) & Decision \\
    \midrule
    \texttt{beta}           & 2 (ML feature)     & \textit{TBD} & \textit{TBD} \\
    \texttt{chi2pid}        & 2 (ML feature)     & \textit{TBD} & \textit{TBD} \\
    \texttt{ftof\_energy\_1A} & 2 (ML feature)   & \textit{TBD} & \textit{TBD} \\
    \texttt{ftof\_energy\_1B} & 2 (ML feature)   & \textit{TBD} & \textit{TBD} \\
    \texttt{ftof\_time\_1A}   & 2 (ML feature)   & \textit{TBD} & \textit{TBD} \\
    \texttt{ftof\_time\_1B}   & 2 (ML feature)   & \textit{TBD} & \textit{TBD} \\
    \texttt{ecin\_energy}   & 2 (ML feature)     & \textit{TBD} & \textit{TBD} \\
    \texttt{ecout\_energy}  & 2 (ML feature)     & \textit{TBD} & \textit{TBD} \\
    \texttt{nphe\_htcc}     & 2 (ML feature)     & \textit{TBD} & \textit{TBD} \\
    \texttt{pcal\_energy}   & 3 (candidate)      & \textit{TBD} & \textit{TBD} \\
    \texttt{ftof\_energy\_2} & 3 (candidate)     & \textit{TBD} & \textit{TBD} \\
    \midrule
    \multicolumn{4}{l}{\textit{(complete with all Group 1--3 variables after audit)}} \\
    \bottomrule
  \end{tabular}
\end{table}

% Cooper: add a paragraph here once the audit is complete.
% TODO: discuss findings.

% =============================================================================
\section{Summary and Next Steps}
\label{sec:summary}
% Cooper: write 2-3 paragraphs covering what was accomplished in Weeks 1-2:
%   (1) ntuple produced and validated (row count, class balance, RICH presence);
%   (2) baseline PID characterized (cite your headline numbers from Table 1
%       and Figure 4 once filled in);
%   (3) data/MC audit complete (cite the feature decision table).
% Then one short paragraph pointing to Weeks 3+ work.

Weeks~1 and~2 established the foundations for the ML classification study.
The \texttt{clasdis} MC ntuple was produced via
\texttt{processing\_mc\_pid\_training.groovy} and loaded into Python;
initial diagnostic plots (Figures~\ref{fig:beta_vs_p}
and~\ref{fig:chi2pid_dist}) confirmed that the truth-class separation problem
is clearly visible in the raw detector variables.
% TODO: add specific numbers (row count, class balance, RICH status) after
% completing Task 3.

The baseline EB pass-2 momentum-dependent \texttt{chi2pid} cut was applied and
characterized via the formal PID metrics of Section~\ref{sec:metrics}.
% TODO: summarize the headline baseline contamination number from Table 1 once
% filled in (e.g., "pion contamination C^{pi->K} reaches X% in the highest
% momentum bin p in [4,5] GeV/c").

The data/MC variable agreement audit (Section~\ref{sec:datamc}) established
which detector variables may enter ML training.
% TODO: summarize the audit outcome: how many variables KEEP/CANDIDATE/DROP.

Weeks~3 through~10 proceed to full-statistics ntuple production, BDT and MLP
model training, per-$(p,\theta)$-bin threshold optimization, and data-driven
validation using the exclusive $ep \to e\,h^+\,(n)$ missing-mass method and
RICH cross-check.
The final deliverable is a calibrated probabilistic classifier with a
quantified contamination improvement over the EB baseline.

% =============================================================================
\appendix
\section{Output Ntuple Column Reference}
\label{app:columns}
% Cooper: verify the column count (53 vs 50 vs 57 — check the script's
% startup banner when you run it) and correct the table below if needed.
% The column map in cooper_day1_and_week1.md Section 4e is the reference.

Table~\ref{tab:columns} lists all 53~columns in the ntuple produced by
\texttt{processing\_mc\_pid\_training.groovy}.

\begin{table}[h]
  \centering
  \caption{Output ntuple column reference (53 columns).
           Sentinel value $-9999$ indicates missing or inapplicable data.}
  \label{tab:columns}
  \small
  \begin{tabular}{clll}
    \toprule
    Col(s) & Name(s) & Group & Description \\
    \midrule
    1--3   & \texttt{runnum, evnum, helicity}          & Event ID      & Run number, event number, beam helicity \\
    4--8   & \texttt{Q2, W, x, y, nu}                  & DIS kin.      & Standard DIS variables \\
    9      & \texttt{pid}                               & Track ID      & EB-assigned PDG code \\
    10--15 & \texttt{p, theta, phi, vz, sector, status} & Track kin.   & Momentum, angles, vertex $z$, sector, status \\
    16--17 & \texttt{beta, chi2pid}                     & EB inputs     & Measured $\beta$; kaon-hypothesis $\chi^2$ \\
    18--23 & \texttt{ftof\_energy\_1A/1B,}              & FTOF 1A/1B    & Energy, time, path length per panel \\
           & \texttt{ftof\_time\_1A/1B,}                &               & \\
           & \texttt{ftof\_path\_1A/1B}                 &               & \\
    24--29 & \texttt{ecin/ecout\_energy,}               & ECAL in/out   & Energy, time, path length per layer \\
           & \texttt{ecin/ecout\_time,}                  &               & \\
           & \texttt{ecin/ecout\_path}                  &               & \\
    30     & \texttt{nphe\_htcc}                        & HTCC          & Number of HTCC photoelectrons \\
    31--36 & \texttt{pcal\_energy/time/path,}           & PCAL + FTOF-2 & Connor group candidate features \\
           & \texttt{ftof\_energy/time/path\_2}         &               & \\
    37--50 & \texttt{rich\_*} (14 variables)            & RICH          & Cross-check only; not ML training features \\
    51     & \texttt{mc\_matching\_pid}                 & MC truth      & PDG code of matched MC particle ($-9999$ = no match) \\
    52     & \texttt{mc\_parent\_pid}                   & MC truth      & PDG code of parent in Lund record \\
    53     & \texttt{mc\_match\_quality}                & MC truth      & Match quality metric (angular distance proxy) \\
    \bottomrule
  \end{tabular}
\end{table}

% =============================================================================
\bibliographystyle{unsrt}
\bibliography{references}

\end{document}
